El circuito entra a la simulación como netlist de SPICE. La fuente lleva las
dos excitaciones al mismo tiempo: \texttt{AC 1} para el barrido de frecuencia y
\texttt{PULSE} para el escalón. \textsc{ngspice} usa la que corresponde a cada
análisis e ignora la otra, así que un solo archivo alcanza para los dos ensayos.

\begin{lstlisting}[language=spice, caption={\texttt{fuentes/filtro-rlc.cir}}, captionpos=b]
* Filtro pasabajos RLC de segundo orden
Vin in 0 AC 1 PULSE(0 1 0 1u 1u 5m 10m)

R1 in a   330
L1 a  b   100m
RL b  out 38
C1 out 0  220n

.end
\end{lstlisting}

\noindent
Xtal no escribe el análisis adentro del netlist: lo pasa por línea de comandos,
así que el mismo circuito se barre en frecuencia o en el tiempo sin editar el
archivo.
